
\documentclass{article}
\usepackage{colortbl}
\usepackage{makecell}
\usepackage{multirow}
\usepackage{supertabular}

\begin{document}

\newcounter{utterance}

\centering \large Interaction Transcript for game `dond', experiment `coop', episode 4 with mistral{-}large{-}3{-}azure.
\vspace{24pt}

{ \footnotesize  \setcounter{utterance}{1}
\setlength{\tabcolsep}{0pt}
\begin{supertabular}{c@{$\;$}|p{.15\linewidth}@{}p{.15\linewidth}p{.15\linewidth}p{.15\linewidth}p{.15\linewidth}p{.15\linewidth}}
    \# & $\;$A & \multicolumn{4}{c}{Game Master} & $\;\:$B\\
    \hline

    \theutterance \stepcounter{utterance}  
    & & \multicolumn{4}{p{0.6\linewidth}}{
        \cellcolor[rgb]{0.9,0.9,0.9}{
            \makecell[{{p{\linewidth}}}]{
                \texttt{\tiny{[P1$\langle$GM]}}
                \texttt{Je speelt een coöperatief onderhandelingsspel waarin je het eens moet worden met een ander speler over hoe je een set items tussen jullie verdeelt.} \\
\\ 
\texttt{Regels:} \\
\texttt{(a) Jij en de ander speler ontvangen een set items. Elk van jullie krijgt ook een geheime waardefunctie, die aangeeft hoeveel je elk type object waardeert.} \\
\texttt{(b) Je wisselt berichten uit met de ander speler om afspraken te maken over wie wat krijgt. Elk kan maximaal 5 berichten sturen of het spel vroegtijdig beëindigen door op elk moment een geheim voorstel in te dienen.} \\
\texttt{(c) Je wordt gevraagd een geheim voorstel in te dienen dat aangeeft welke items je wilt, opgemaakt tussen vierkante haken als volgt: "{[}Voorstel: <aantal> <objectnaam>, <aantal> <objectnaam>, <...>{]}"} \\
\texttt{(d) Als je voorstellen complementair zijn, dat wil zeggen dat er voldoende items zijn om beide voorstellen in te vullen, krijgt elke speler een score op basis van de som van de waarden van de items die hij heeft ontvangen. Anders krijgen jullie beiden nul punten.} \\
\texttt{(e) Je doel is om de som van je score en die van de ander speler te maximaliseren.} \\
\\ 
\texttt{Laten we beginnen.} \\
\\ 
\texttt{De set beschikbare items is:} \\
\\ 
\texttt{3 boeken, 1 zak, 1 bout} \\
\\ 
\texttt{Je geheime waardefunctie is:} \\
\\ 
\texttt{boek: 0, zak: 2, bout: 8} \\
\\ 
\texttt{BELANGRIJK: Je berichten, tenzij het een geheim voorstel is, worden rechtstreeks aan de ander speler doorgegeven, dus voeg geen antwoord op de regels of tekst toe die je bericht aankondigt. Om een geheim voorstel in te dienen, gebruik je het aangegeven format. Gebruik geen vierkante haken wanneer je met de ander speler communiceert, anders wordt dit geïnterpreteerd als je geheime voorstel.} \\
            }
        }
    }
    & \\ \\

    \theutterance \stepcounter{utterance}  
    & \multicolumn{4}{p{0.6\linewidth}}{
        \cellcolor[rgb]{0.9,0.9,0.9}{
            \makecell[{{p{\linewidth}}}]{
                \texttt{\tiny{[P1$\rangle$GM]}}
                \texttt{Ik stel voor dat jij de bout neemt en ik de zak, aangezien de bout voor jou het meest waard is en de zak voor mij het op één na meest waardevol is. Wat denk je hiervan?} \\
            }
        }
    }
    & & \\ \\

    \theutterance \stepcounter{utterance}  
    & & & \multicolumn{2}{p{0.3\linewidth}}{
        \cellcolor[rgb]{0.9,0.9,0.9}{
            \makecell[{{p{\linewidth}}}]{
                \texttt{\tiny{[GM$|$GM]}}
                \texttt{continue} \\
            }
        }
    }
    & & \\ \\

    \theutterance \stepcounter{utterance}  
    & & & \multicolumn{2}{p{0.3\linewidth}}{
        \cellcolor[rgb]{0.9,0.9,0.9}{
            \makecell[{{p{\linewidth}}}]{
                \texttt{\tiny{[GM$|$GM]}}
                \texttt{Ik stel voor dat jij de bout neemt en ik de zak, aangezien de bout voor jou het meest waard is en de zak voor mij het op één na meest waardevol is. Wat denk je hiervan?} \\
            }
        }
    }
    & & \\ \\

    \theutterance \stepcounter{utterance}  
    & & \multicolumn{4}{p{0.6\linewidth}}{
        \cellcolor[rgb]{0.9,0.9,0.9}{
            \makecell[{{p{\linewidth}}}]{
                \texttt{\tiny{[GM$\rangle$P2]}}
                \texttt{Je speelt een coöperatief onderhandelingsspel waarin je het eens moet worden met een ander speler over hoe je een set items tussen jullie verdeelt.} \\
\\ 
\texttt{Regels:} \\
\texttt{(a) Jij en de ander speler ontvangen een set items. Elk van jullie krijgt ook een geheime waardefunctie, die aangeeft hoeveel je elk type object waardeert.} \\
\texttt{(b) Je wisselt berichten uit met de ander speler om afspraken te maken over wie wat krijgt. Elk kan maximaal 5 berichten sturen of het spel vroegtijdig beëindigen door op elk moment een geheim voorstel in te dienen.} \\
\texttt{(c) Je wordt gevraagd een geheim voorstel in te dienen dat aangeeft welke items je wilt, opgemaakt tussen vierkante haken als volgt: "{[}Voorstel: <aantal> <objectnaam>, <aantal> <objectnaam>, <...>{]}"} \\
\texttt{(d) Als je voorstellen complementair zijn, dat wil zeggen dat er voldoende items zijn om beide voorstellen in te vullen, krijgt elke speler een score op basis van de som van de waarden van de items die hij heeft ontvangen. Anders krijgen jullie beiden nul punten.} \\
\texttt{(e) Je doel is om de som van je score en die van de ander speler te maximaliseren.} \\
\\ 
\texttt{Laten we beginnen.} \\
\\ 
\texttt{De set beschikbare items is:} \\
\\ 
\texttt{3 boeken, 1 zak, 1 bout} \\
\\ 
\texttt{Je geheime waardefunctie is:} \\
\\ 
\texttt{boek: 1, zak: 2, bout: 5} \\
\\ 
\texttt{BELANGRIJK: Je berichten, tenzij het een geheim voorstel is, worden rechtstreeks aan de ander speler doorgegeven, dus voeg geen antwoord op de regels of tekst toe die je bericht aankondigt. Om een geheim voorstel in te dienen, gebruik je het aangegeven format. Gebruik geen vierkante haken wanneer je met de ander speler communiceert, anders wordt dit geïnterpreteerd als je geheime voorstel.} \\
\\ 
\texttt{Ik stel voor dat jij de bout neemt en ik de zak, aangezien de bout voor jou het meest waard is en de zak voor mij het op één na meest waardevol is. Wat denk je hiervan?} \\
            }
        }
    }
    & \\ \\

    \theutterance \stepcounter{utterance}  
    & & & \multicolumn{4}{p{0.6\linewidth}}{
        \cellcolor[rgb]{0.9,0.9,0.9}{
            \makecell[{{p{\linewidth}}}]{
                \texttt{\tiny{[GM$\langle$P2]}}
                \texttt{Ik stel voor dat we de verdeling als volgt doen: ik neem 2 boeken en jij neemt de bout, de zak en het overgebleven boek. Wat vind je daarvan?} \\
            }
        }
    }
     \\ \\

    \theutterance \stepcounter{utterance}  
    & & & \multicolumn{2}{p{0.3\linewidth}}{
        \cellcolor[rgb]{0.9,0.9,0.9}{
            \makecell[{{p{\linewidth}}}]{
                \texttt{\tiny{[GM$|$GM]}}
                \texttt{continue} \\
            }
        }
    }
    & & \\ \\

    \theutterance \stepcounter{utterance}  
    & & & \multicolumn{2}{p{0.3\linewidth}}{
        \cellcolor[rgb]{0.9,0.9,0.9}{
            \makecell[{{p{\linewidth}}}]{
                \texttt{\tiny{[GM$|$GM]}}
                \texttt{Ik stel voor dat we de verdeling als volgt doen: ik neem 2 boeken en jij neemt de bout, de zak en het overgebleven boek. Wat vind je daarvan?} \\
            }
        }
    }
    & & \\ \\

    \theutterance \stepcounter{utterance}  
    & & \multicolumn{4}{p{0.6\linewidth}}{
        \cellcolor[rgb]{0.9,0.9,0.9}{
            \makecell[{{p{\linewidth}}}]{
                \texttt{\tiny{[P1$\langle$GM]}}
                \texttt{Ik stel voor dat we de verdeling als volgt doen: ik neem 2 boeken en jij neemt de bout, de zak en het overgebleven boek. Wat vind je daarvan?} \\
            }
        }
    }
    & \\ \\

    \theutterance \stepcounter{utterance}  
    & \multicolumn{4}{p{0.6\linewidth}}{
        \cellcolor[rgb]{0.9,0.9,0.9}{
            \makecell[{{p{\linewidth}}}]{
                \texttt{\tiny{[P1$\rangle$GM]}}
                \texttt{Dat voorstel geeft mij een waarde van 10 (bout + zak) en jou 0 (boeken), wat niet eerlijk verdeeld is gezien mijn waardefunctie. Laten we een verdeling zoeken waarbij we allebei meer waarde krijgen.} \\
\\ 
\texttt{Ik ben vooral geïnteresseerd in de bout (8) en de zak (2). Jij lijkt vooral waarde te hechten aan de boeken. Zou je akkoord gaan met: jij krijgt alle 3 boeken, en ik krijg de bout en de zak? Zo krijgen we allebei maximale waarde uit onze eigen voorkeuren.} \\
            }
        }
    }
    & & \\ \\

    \theutterance \stepcounter{utterance}  
    & & & \multicolumn{2}{p{0.3\linewidth}}{
        \cellcolor[rgb]{0.9,0.9,0.9}{
            \makecell[{{p{\linewidth}}}]{
                \texttt{\tiny{[GM$|$GM]}}
                \texttt{continue} \\
            }
        }
    }
    & & \\ \\

    \theutterance \stepcounter{utterance}  
    & & & \multicolumn{2}{p{0.3\linewidth}}{
        \cellcolor[rgb]{0.9,0.9,0.9}{
            \makecell[{{p{\linewidth}}}]{
                \texttt{\tiny{[GM$|$GM]}}
                \texttt{Dat voorstel geeft mij een waarde van 10 (bout + zak) en jou 0 (boeken), wat niet eerlijk verdeeld is gezien mijn waardefunctie. Laten we een verdeling zoeken waarbij we allebei meer waarde krijgen.} \\
\\ 
\texttt{Ik ben vooral geïnteresseerd in de bout (8) en de zak (2). Jij lijkt vooral waarde te hechten aan de boeken. Zou je akkoord gaan met: jij krijgt alle 3 boeken, en ik krijg de bout en de zak? Zo krijgen we allebei maximale waarde uit onze eigen voorkeuren.} \\
            }
        }
    }
    & & \\ \\

    \theutterance \stepcounter{utterance}  
    & & \multicolumn{4}{p{0.6\linewidth}}{
        \cellcolor[rgb]{0.9,0.9,0.9}{
            \makecell[{{p{\linewidth}}}]{
                \texttt{\tiny{[GM$\rangle$P2]}}
                \texttt{Dat voorstel geeft mij een waarde van 10 (bout + zak) en jou 0 (boeken), wat niet eerlijk verdeeld is gezien mijn waardefunctie. Laten we een verdeling zoeken waarbij we allebei meer waarde krijgen.} \\
\\ 
\texttt{Ik ben vooral geïnteresseerd in de bout (8) en de zak (2). Jij lijkt vooral waarde te hechten aan de boeken. Zou je akkoord gaan met: jij krijgt alle 3 boeken, en ik krijg de bout en de zak? Zo krijgen we allebei maximale waarde uit onze eigen voorkeuren.} \\
            }
        }
    }
    & \\ \\

    \theutterance \stepcounter{utterance}  
    & & & \multicolumn{4}{p{0.6\linewidth}}{
        \cellcolor[rgb]{0.9,0.9,0.9}{
            \makecell[{{p{\linewidth}}}]{
                \texttt{\tiny{[GM$\langle$P2]}}
                \texttt{{[}Voorstel: 3 boeken{]}} \\
            }
        }
    }
     \\ \\

    \theutterance \stepcounter{utterance}  
    & & & \multicolumn{2}{p{0.3\linewidth}}{
        \cellcolor[rgb]{0.9,0.9,0.9}{
            \makecell[{{p{\linewidth}}}]{
                \texttt{\tiny{[GM$|$GM]}}
                \texttt{proposal} \\
            }
        }
    }
    & & \\ \\

    \theutterance \stepcounter{utterance}  
    & & & \multicolumn{2}{p{0.3\linewidth}}{
        \cellcolor[rgb]{0.9,0.9,0.9}{
            \makecell[{{p{\linewidth}}}]{
                \texttt{\tiny{[GM$|$GM]}}
                [3, 0, 0]
            }
        }
    }
    & & \\ \\

    \theutterance \stepcounter{utterance}  
    & & \multicolumn{4}{p{0.6\linewidth}}{
        \cellcolor[rgb]{0.9,0.9,0.9}{
            \makecell[{{p{\linewidth}}}]{
                \texttt{\tiny{[P1$\langle$GM]}}
                \texttt{De andere speler heeft een geheim voorstel ingediend. Het is nu jouw beurt om een eigen voorstel in te dienen.} \\
            }
        }
    }
    & \\ \\

    \theutterance \stepcounter{utterance}  
    & \multicolumn{4}{p{0.6\linewidth}}{
        \cellcolor[rgb]{0.9,0.9,0.9}{
            \makecell[{{p{\linewidth}}}]{
                \texttt{\tiny{[P1$\rangle$GM]}}
                \texttt{{[}Voorstel: 1 bout, 1 zak{]}} \\
            }
        }
    }
    & & \\ \\

    \theutterance \stepcounter{utterance}  
    & & & \multicolumn{2}{p{0.3\linewidth}}{
        \cellcolor[rgb]{0.9,0.9,0.9}{
            \makecell[{{p{\linewidth}}}]{
                \texttt{\tiny{[GM$|$GM]}}
                \texttt{proposal} \\
            }
        }
    }
    & & \\ \\

    \theutterance \stepcounter{utterance}  
    & & & \multicolumn{2}{p{0.3\linewidth}}{
        \cellcolor[rgb]{0.9,0.9,0.9}{
            \makecell[{{p{\linewidth}}}]{
                \texttt{\tiny{[GM$|$GM]}}
                [0, 1, 1]
            }
        }
    }
    & & \\ \\

    \theutterance \stepcounter{utterance}  
    & & & \multicolumn{2}{p{0.3\linewidth}}{
        \cellcolor[rgb]{0.9,0.9,0.9}{
            \makecell[{{p{\linewidth}}}]{
                \texttt{\tiny{[GM$|$GM]}}
                [[0, 1, 1], [3, 0, 0]]
            }
        }
    }
    & & \\ \\

\end{supertabular}
}

\end{document}
